
\section{Introduction} \label{sec:introduction}
We consider the denoising problem, where the goal is to remove noise from an unknown image or signal.
In more detail, our goal is to obtain an estimate of an image or signal $y_\ast \in \R^n$ from a noisy measurement
\[
y = y_\ast + \noise.
\]
Here, $\noise$ is unknown noise, which we model as a zero-mean white Gaussian random variable with covariance matrix $\sigma^2/n I$.
Image denoising relies on generative or prior assumptions on the image $y_\ast$, 
such as self-similarity within images~\citep{dabov_image_2007},
 sparsity in fixed~\citep{donoho_de-noising_1995} and learned bases~\citep{elad_image_2006}, and most recently, by assuming the image can be generated by a pre-trained deep-neural network~\citep{burger_image_2012,zhang_beyond_2017}.
Deep-network based approaches, typically yield the best denoising performance. This success can be attributed to their ability to efficiently represent and learn realistic image priors, for example via auto-decoders~\citep{hinton_reducing_2006} and generative adversarial models~\citep{goodfellow_generative_2014}.
%Over the last few years, the quality of deep priors has significantly improved \citep{Karras17, DIP}.  As this field matures, priors will be developed with even smaller latent code dimensionality and more accurate approximation of natural signal manifolds. Consequently, the representation error from deep priors will decrease, and thereby enable even more powerful denoisers.


%Motivated by the fact that the latter, deep learning based denoisers achieve state-of-the-art results, 
Motivated by this success story, we assume that the image $y_\ast$ lies in the range of an image-generating network. 
In this paper, we propose the first algorithm for solving denoising with deep generative priors that provably finds an approximation of the underlying image. 
%Following the lead of \cite{hand_global_2017}, we assume an expansive Gaussian model for the deep generative network in order to establish this result. 
As the influence of deep networks in denoising and inverse problems grows, it becomes increasingly important to understand their performance at a theoretical level.  Given that most optimization approaches for deep learning are first-order gradient methods, a justification is needed for why they do not get stuck in local minima.  

The most related work that establishes theoretical reasons for why gradient methods might succeed when using deep generative priors for solving inverse problems, is \citep{hand_global_2017}.  In it, the authors establish global favorability for optimization of a $\ell_2$-loss function under a random neural network model.  Specifically, they show existence of a descent direction outside a ball around the global optimizer and a negative multiple of it in the latent space of the generative model.  This work does not justify why the one spurious point is avoided by gradient descent, nor does it provide a specific algorithm which provably estimates the global minimizer, nor does it provide an analysis of the robustness of the problem with respect to noise.  This work was subsequently extended to include the case of generative convolutional neural networks by \cite{ma2018invertibility}, but that work too does not prove convergence of a specific algorithm.

% mapping a low-dimensional latent space to a high-dimensional image.
%$G \colon \reals^k \to \reals^n$. Here, the $k$-dimensional input to the network can be viewed as a latent, low-dimensional code for the high dimensional image. We estimate the

%  While this assumption on the network is strong, it is the weakest known condition that permits such a claim.  Further, empirical evidence from trained networks show that they can have statistics consistent with Gaussians \citep{Arora15}, and truly random nets, such as in the Deep Image Prior \citep{DIP}, are increasingly becoming of practical relevance.  Thus, theoretical advances on random nets is of independent interest.  While trained networks are not exactly Gaussian, it is not clear which non-Gaussian distributions of network weights would be considered sufficiently realistic.   The work in \cite{hand_global_2017} leaves open the question about what specific algorithm will actually find an approximation of the global optimizer.  Further, it entirely leaves open the question of how noise affects the empirical risk optimization.  

%While deep networks can model complex priors through many parameters and non-linearities, those non-linearities also make their analysis inherently difficult, increasing the gap between theory and practice. As a consequence, there is a need for theory explaining the success of deep network based priors.

\paragraph{Contributions:}
The goal of this paper is to analytically quantify the denoising performance of deep-prior based denoisers.
% We propose a simple and efficient algorithm for denoising based on a based on a $d$-layer generative neural network $G\colon \R^k \to \R^{n}$, with $k<n$, and quantify its denoising performance for
Specifically, we characterize the performance of two simple and efficient algorithms for denoising based on a $d$-layer generative neural network
$G\colon \R^k \to \R^{n}$, with $k<n$.

We first provide a simple result for an encoder-generator network $G( E(y) )$ where $E \colon \R^n \to \R^k$ is an encoder network.
We show that if we pass noise through an encoder-decoder network $G( E(y) )$ that acts as the identity on a class of images of interest, then it reduces the random noise by $O(k/n)$. 
This result requires no assumptions on the weights of the network.

The second and main result of our paper pertains to denoising by optimizing over the latent code of a generator network with random weights.
We propose a gradient method 
%with a tweak 
that attempts to minimize the least-squares loss $f(x) = \frac{1}{2} \norm{G(x) - y}^2$ between the noisy image $y$ and an image in the range of the generator, $G(x)$. Even though $f$ is non-convex, we show that a gradient method yields an estimate $\hat x$ obeying
\begin{align*}
\norm{G(\hat x) - y_\ast}^2
\lesssim \sigma^2 \frac{k}{n},
\end{align*}
with high probability, where the notation $\lesssim$ absorbs a constant factor depending on the number of layers of the network, and its expansitivity, as discussed in more detail later.
Our result shows that the denoising rate of a deep generator based denoiser is optimal in terms of the dimension of the latent representation.
We also show in numerical experiments, that this rate---shown to be analytically achieved for random priors---is also experimentally achieved for priors learned from real imaging data.
%Up to a logarithmic factor, this result sharply characterizes the denoising rate of generative neural networks.

\paragraph{Related work:}
We hasten to add that a close theoretical work to the question considered in this paper is the paper~\citep{bora_compressed_2017}, which solves a noisy compressive sensing problem with generative priors by minimizing an $\ell_2$-loss. Under the assumption that the network is Lipschitz, they show that if the global optimizer can be found, which is in principle NP-hard, then a signal estimate is recovered to within the noise level.  While the Lipschitzness assumption is quite mild, the resulting theory does not provide justification for why global optimality can be reached.

%%%

\section{Background on denoising with classical and deep-learning based methods}

As mentioned before, image denoising relies on modeling or prior assumptions on the image $y_\ast$. For example, suppose that the image $y_\ast$ lies in a $k$-dimensional subspace of $\mathbb R^n$ denoted by $\mathcal Y$.
Then we can estimate the original image by finding the closest point in $\ell_2$-distance to the noisy observation $y$ on the subspace $\mathcal Y$.
The corresponding estimate, denoted by $\hat y$, obeys
%$\hat y = \arg \min_{y'\in \mathcal Y} \norm{y' - y}^2$
\begin{align}
\label{eq:subspacerate}
\norm{\hat y - y_\ast}^2 \lesssim \sigma^2 \frac{k}{n},
\end{align}
with high probability (throughout, $\norm{\cdot}$ denotes the $\ell_2$-norm).
Thus, the noise energy is reduced by a factor of $k/n$ over the trivial estimate $\hat y = y$ which does not use any prior knowledge of the signal.
The denoising rate~\eqref{eq:subspacerate} shows that the more concise the image prior or image representation (i.e., the smaller $k$), the more noise can be removed. If on the other hand the image model (the subspace, in this example) does not include the original image $y_\ast$, then the error bound~\eqref{eq:subspacerate} increases, as we would remove a significant part of the signal along with noise when projecting onto the range of the image prior.
Thus a concise and accurate model is crucial for denoising.

Real world signals rarely lie in \emph{a priori} known subspaces, and the last few decades of image denoising research have developed sophisticated algorithms based on accurate image models.
Examples include algorithms based on sparse representations in overcomplete dictionaries such as wavelets~\citep{donoho_de-noising_1995} and curvelets~\citep{starck_curvelet_2002}, and algorithms based on exploiting self-similarity within images~\citep{dabov_image_2007}.
A prominent example of the former class of algorithms is the (state-of-the-art) BM3D~\citep{dabov_image_2007} algorithm. %, which achieves state-of-the-art performance for certain denoising problems.
However, the nuances of real world images are difficult to describe with handcrafted models. Thus, starting with the paper~\citep{elad_image_2006} that proposes to learn sparse representation based on training data,
it has become common to learn concise representation for denoising (and other inverse problems) from a set of training images.

\citet{burger_image_2012} applied deep networks to the denoising problem, by training a plain neural network on a large set of images.
Since then, deep learning based denoisers~\citep{zhang_beyond_2017} have set the standard for denoising. The success of deep network priors can be attributed to their ability to efficiently represent and learn realistic image priors, for example via auto-decoders~\citep{hinton_reducing_2006} and generative adversarial models~\citep{goodfellow_generative_2014}.
Over the last few years, the quality of deep priors has significantly improved \citep{Karras17, DIP,heckel_deep_2019}. As this field matures, priors will be developed with even smaller latent code dimensionality and more accurate approximation of natural signal manifolds. Consequently, the representation error from deep priors will decrease, and thereby enable even more powerful denoisers.

%%%

\section{Denoising with a neural network with an hourglass structure}

Perhaps the most straight-forward and classical approach to using deep networks for denoising is to train a deep network with an autoencoder or hourglass structure end-to-end to perform denoising.   
%Traditionally, autoencoders have been used to learn feature representations for some data set.
%Denoising autoencoders artificially corrupt the input data in order to force the network to learn a more robust representation of the data. 
%In either case, an autoencoder is a neural network mapping an image to an image ($G \colon \reals^n \to \reals^n$) which has the property such that $G(x) \approx x$ for $x$ in a data set of interest. An autoencoder compresses data from the input layer into a short code, and then generates an output from the short code. 
%The length of the short code determines the denoising performance.
An autoencoder compresses data from the input layer into a low-dimensional code, and then generates an output from that code. In this section, we analyze such networks from the perspective of denoising.

Specifically, we show mathematically that a simple model for neural networks with an hourglass structure achieves optimal denoising rates, as given by the dimensionality of the low-dimensional code.
%We next provide an illustrative example showing that the size of the short code determines the amount of noise reduction or the denoising rate.
An autoencoder $H(x) = G( E(x) )$ 
%maps an image to an image, and 
consists of an encoder network $E \colon \reals^n \to \reals^k$ mapping an image to a low-dimensional latent representation, and a decoder or generator network 
$G \colon \reals^k \to \reals^n$
mapping the latent representation to an image. 
To see that the size of the low-dimensional code, $k$, determines the denoising rate, consider a simple one-layer encoder and multilayer decoder of the form
\begin{align}
\label{eq:encoderdecoder}
E(y)
=
\relu(W'y)
,
\quad
G(x) = \relu(\Win{d} \ldots \relu (\Win{2} \relu(\Win{1} x)) \ldots ),
\end{align}
where $\relu(x) = \max(x, 0)$ applies entrywise,
$W'$ are the weights of the encoder, 
$W_i \in \R^{n_i \times n_{i-1}}$, are the weights in the $i$-th layer of the decoder, and $n_i$ is the number of neurons in the $i$th layer.

Typically, the autoencoder network $H$ is trained such that 
$H(y) \approx y$ for some class of signals of interest (say, natural images).
The following proposition shows that the structure of the network alone guarantees that an autoencoder ``filters out'' most of the noise. %which can be formalized with the following  proposition that is based on a subspace counting and noise concentration argument.

\begin{proposition}
\label{prop:autodec}
Let $H = G \circ E$ be an autoencoder of the form~\eqref{eq:encoderdecoder} and note that it is piecewise linear, i.e., $H(y) = Uy$ in a region around $y$. 
Suppose that $\norm{U}^2 \leq  2$ for all such  regions. 
Let $\noise$ be Gaussian noise with covariance matrix $\sigma I$, $\sigma > 0$. Then, provided that 
$ k\cdot 32 \log(2n_1 n_2 \ldots n_d)   \leq n$, we have 
with probability at least $1 - 2 e^{-k \log(2n_1 n_2 \ldots n_d)}$, that 
\[
\norm[2]{H(\eta)}^2 \leq  
5\frac{k}{n} \log(2n_1 n_2 \ldots n_d) \norm[2]{\eta}^2.
\]
\end{proposition}

Note that the assumption $\norm{U}^2 \leq 2$ implies that
$\norm[2]{H(y)}^2 / \norm[2]{y}^2 \leq 2$ for all $y$. This in turn guarantees that the autoencoder does not ``amplify'' a signal too much.
 Note that we envision an autoencoder that is trained such that it obeys $H(y) \approx y$ for $y$ in a class of images. %Such networks would approximately satisfy the property that  $\norm[2]{H(y)}^2 / \norm[2]{y}^2 \leq 2$.  
 The proposition would then justify why the the feedforward network reduces noise by $O(k/n)$. %Thus, the amount of noise that an autodecoder does not filter out is determined by the size of the latent code, $k$.

The proof of Proposition~\ref{prop:autodec}, contained in the appendix, shows that $H$ lies in the range of a union of $k$-dimensional subspaces and then uses a standard concentration argument showing that the union of those subspaces can represent no more than a fraction of $O(k/n)$ of the noise.

In the remainder of the paper we shows that denoising via enforcing a generative prior gives us an analogous denoising rate.


